% =============================================================================
% Section 6: The finite-local audit chain
% Negative-program backbone. Highest-risk for overclaiming — kept tight on scope.
% =============================================================================

\section{The finite-local audit chain}
\label{sec:csp-audit}

This section establishes the negative content of the paper. We
characterize the family of \emph{finite-local forced-divisor
certificate primitives} we have tested against the residual frontier
of \cref{thm:143-collapse}, summarize the experimental and structural
results across each layer, and conclude with a precise statement of
what the audit chain does and does not rule out.

\subsection{What counts as a finite-local certificate}
\label{sec:ple-formal}

\begin{definition}[Finite-local forced-divisor certificate]
\label{def:finite-local-cert}
Fix a residue class $r \in \Ropen$. A \emph{finite-local
forced-divisor certificate} for $r$ is a finite-data proof that for
every integer $n \equiv r \pmod M$ with $n > n_0$ (an explicit
threshold), there exists $1 \le k \le K$ with
\[
\tauf\bigl( n - k \bigr) \;>\; k + 2 ,
\]
established from local valuation data at a finite set of primes
$P$. The proof must be combinatorial and exhibit, for each $n$ in the
class, an explicit such $k$ together with a divisor $d \mid (n - k)$
with $d \le \prod_{p \in P} p^{E_p}$ and $\tauf(d) > k + 2$.
\end{definition}

The class \PLEclass\ of \cref{def:ple} is the natural completion of
\cref{def:finite-local-cert} under all primitives that can be
expressed as decidable Boolean combinations of local valuation
constraints. The audit chain below tests \PLEclass\ at increasing
expressive power.

\subsection{Cycle-3 valuation leaves}
\label{sec:cycle3-audit}

The earliest closure primitive in the project, the
\emph{cycle-3 valuation leaf}, attempts to certify that for some
small fixed $k$ the forced divisor mass at three specific primes
already exceeds $k + 2$. For
$k \le K_0 = 95$ the available budget $k + 2$ allows at most
$\tauf(D_k) \le (k+2)/2$ on the kernel-forced divisor (Conjecture B
form), and the cycle-3 enumeration shows no triple of primes
$(p, q, r)$ in the residue's structure suffices to exceed the budget
at any of the protected shifts. The cycle-3 framework is therefore
\emph{insufficient} for closure: it produces no certificate.

\subsection{DivisorMask levels 1--3}
\label{sec:divisormask-audit}

The DivisorMask family generalizes cycle-3 by allowing arbitrary
prime support for the forced divisor. We tested:

\begin{itemize}[leftmargin=*]
\item \textbf{L1 (squarefree free-prime DivisorMask):} forced divisor
  of squarefree shape, primes ranging in a free set
  $P \setminus \{2, 3, 5, 7\}$.
\item \textbf{L2 (joint $p$-adic DivisorMask):} arbitrary $p$-adic
  valuations $v_p$ on a fixed prime set, with combinatorial
  constraints across shifts.
\item \textbf{L3 (exact-valuation $\mu$-DivisorMask):} exact (rather
  than lower-bound) valuations, with $\mu$-tracking of free vs.\
  forced contributions.
\end{itemize}

Each of L1--L3 was encoded as a CP-SAT instance with the
\PLEclass\ structure of \cref{def:ple} and tested on representative
residues from $\Ropen$, including the sentinels $r \in \{1716, 4862\}$
where the strongest finite-local pressure was observed. \emph{Each
level yielded SAT escapes at the parameters tested}; no closure
certificate was produced. The escape patterns at L1--L3 informed
the design of the prime-residual kernel framework below.

\subsection{Prime-residual and Hensel-balanced kernels}
\label{sec:hensel-kernel-audit}

The prime-residual kernel framework (Levels~4--6) introduces an
algebraic structure rather than enumerating divisor patterns
combinatorially. For each residue $r$ and each fixed window length
$K_0$, the Hensel-complete kernel construction produces a CRT class
$s \equiv a \pmod{\Lcap}$ such that
\[
F_k(s) \;=\; D_k(a) \cdot Q_k(t), \qquad
Q_k(t) = \alpha_k + \beta_k \cdot t, \quad \beta_k > 0,
\]
for every $1 \le k \le K_0$, with $D_k(a) \mid \Acap \cdot \Lcap$
and the admissibility condition $2 \tauf(D_k(a)) \le k + 2$
(\emph{Conjecture B}). The kernel verifier (cf.\
\cref{sec:lean}) confirms admissibility via six independent checks
including $p$-adic admissibility for every $p \le K_0$ and structural
admissibility for $p > K_0$.

We summarize the empirical state through $K_0 \in \{95, 200, 1021\}$:
\emph{Hensel-complete kernels exist for every residue
$r \in \Ropen$ at every tested $K_0$}, with the prime-residual
factorization and Conjecture B both holding throughout. This is the
\emph{barrier}: rather than ruling out a counterexample, the kernel
framework exhibits a survivor-compatible algebraic structure for
every shift in the protected window.

\begin{remark}[Class-specific structural certificates do not extend]
\label{rem:class-specific-struct-certs}
At parameters $K_{\rm scan} \le 10^7$ a structural-kill detector
(an internal audit instrument verifying that named residues fail
under the prime-tuple admissibility constraints) identifies four
residues with kills: $r \in \{858, 1716, 2431, 4862\}$, the latter two at
$k = 132$ and $k = 280$ respectively (\emph{small-shift sparse-prime
kills}), and the former two at $k = 2520 r$ (\emph{mega-shift dense-prime
artifacts}; $F_k(s) = \Acap \cdot s$ when $c = 0$). \emph{None of
these certificates extend to a base-AP-universal closure}; the
kernel-space CSP exhibits escaping kernel bases for the small-shift
cases, and the mega-shift artifacts apply only to the JSON kernel
CRT class chosen by the solver, not to all $n$ in the residue class
(per a separate quantifier audit confirming no $\forall \exists$ inversions in the closure proof).
\end{remark}

\subsection{Moving-window interval-cover audit}
\label{sec:interval-cover-audit}

The strongest finite-local primitive we tested combines the kernel
framework above with auxiliary safe-set residues at primes beyond
the kernel cap structure.

\begin{definition}[Full safe-set CSP]
\label{def:full-safe-set-csp}
Fix $r \in \Ropen$, a kernel $(a, \Lcap)$ for $r$, a protected
window $K_0$, an auxiliary prime bound $B$, and a scan bound
$K_{\rm scan}$. The \emph{full safe-set CSP} $\Phi_r(K_0, B,
K_{\rm scan})$ has variables:
\begin{itemize}[leftmargin=*]
\item $a_p \in \Z / p^{E_p}\Z$ for each prime $p$ in the kernel
  cap structure, varying the kernel CRT base subject to Hensel
  admissibility;
\item $y_\ell \in S_\ell \subset \Z / \ell\Z$ for each auxiliary
  prime $\ell$ with $89 < \ell \le B$, where $S_\ell$ is the safe
  set of \cref{sec:setup}.
\end{itemize}
The CSP is \emph{satisfiable} iff there is a joint assignment such
that for every shift $k$ in $(K_0, K_{\rm scan}]$, the forced
divisor $D_k$ extracted from the local data does not satisfy
$\tauf(D_k) > k + 2$.
\end{definition}

\begin{theorem}[Interval-cover insufficiency; SAT-certified]
\label{thm:interval-cover-insufficiency}
For sentinel residues $r \in \{1716, 4862\}$ and parameters
$(K_0, B, K_{\rm scan})$ with $K_0 \in \{95, 200, 1021\}$,
$B \le 1000$, $K_{\rm scan} \le 10^5$, the full safe-set CSP
$\Phi_r(K_0, B, K_{\rm scan})$ of
\cref{def:full-safe-set-csp} is satisfiable. The satisfiability is
witnessed by an explicit assignment, certified by a CaDiCaL-generated
LRAT proof, and verified by the \texttt{cake\_lpr} CakeML checker.
\end{theorem}

The certificate files and verifier output are described in
\cref{sec:lean}. We emphasize that the same audit was extended to a
\emph{nested} formulation in which the protected and scan windows
overlap (verified by direct enumeration of the nested-survivor tree at depth $1021$); the nested CSP
collapses to the depth-$1021$ CSP for every kernel structure that
satisfies $v_p(F_k(a)) \le v_p(\Acap \cdot \Lcap)$ at all kernel
primes, and produces identical SAT escapes at the same iteration
counts.

\subsection{Audit summary table}
\label{sec:audit-summary}

\begin{table}[h]
\centering
\begin{tabular}{lllc}
\toprule
\textbf{Layer} & \textbf{Primitive family} & \textbf{Result} & \textbf{Status} \\
\midrule
cycle-3 & valuation leaves at three primes
        & forced budget below $k+2$
        & insufficient \\
L1 & squarefree free-prime DivisorMask
        & SAT escapes
        & insufficient \\
L2 & joint $p$-adic DivisorMask
        & SAT escapes
        & insufficient \\
L3 & exact-valuation $\mu$-DivisorMask
        & SAT escapes
        & insufficient \\
L4--L6 & Hensel prime-residual kernels
        & survivor-compatible
        & barrier \\
Hensel $K=10000$ extension
        & constructive smart-lift witnesses
        & $41 / 41$ open-frontier residues satisfy
          $2\tauf(D_k(a)) \le k+2$ for all $k \le 10000$
        & barrier (not closure) \\
struct-kill & fixed-kernel structural shift
        & 4 per-kernel-class certs
        & not base-AP \\
kernel-space & kernel-base CSP
        & SAT escapes at $K_0 \le 1021$
        & closed-negative \\
full safe-set & kernel + aux $S_\ell$ CSP
        & SAT escapes at $B \le 1000$
        & closed-negative \\
nested CSP & overlap protected + scan
        & equivalent to depth-$1021$
        & no new info \\
\midrule
diagonal mega-shift audit
        & identity $F_{r,2520r}(s) = A\cdot s$
        & selected kernel branches for four sentinels are killed at $k_\Delta = 2520r$
        & class-specific, not base-AP \\
branch-gluing Phase B (B1)
        & nested survivor compatibility with diagonal constraint
        & restricted B1 witnesses for $r \in \{1716, 4862\}$ survive protected, diagonal, and scan to $10^7$
        & negative for the tested mechanism \\
diagonal-B1 all-41 triage
        & B1 necessary mode + diagonal at $k_\Delta = 2520\,r$, no scan
        & $40 / 41$ SAT, $1 / 41$ ($r = 0$, $k_\Delta = 0$) is a degenerate edge case
        & closure-incapable on all 41 open residues \\
composite-cover audit
        & composite / non-prime-power forced-divisor clauses
        & $41 / 41$ SAT at $P_{\max} = 199, K_{\max} = 5000$
        & closure-incapable on all 41 open residues \\
small-$k$ quotient-shape compatibility
        & abstract shape of $Q_k = F_k / D_k$ subject to
          multi-shift compatibility $F_i - F_j = j - i$
        & $41 / 41$ SAT for all $K_{\rm small} \le 1021$ via fresh-prime
          assignment $Q_k =$ distinct fresh prime per $k$
        & subsumed by Hensel admissibility ($2\tau(D_k) \le k+2$) \\
\bottomrule
\end{tabular}
\caption{Layers of the finite-local audit chain. ``Insufficient''
denotes a primitive whose forced budget cannot exceed $k+2$ on the
residual frontier; ``barrier'' denotes a primitive that exhibits a
survivor-compatible structure rather than a closure certificate;
``closed-negative'' denotes a SAT-verified failure of the primitive
to close the frontier within the tested caps; ``not base-AP'' and
``class-specific'' denote kernel-class certificates that do not lift
to base-AP closure; ``negative for the tested mechanism'' denotes a
SAT escape under restricted B1 (which proves SAT for the unrestricted
problem) and a witness that survives the structural scan to $K_{\rm scan} = 10^7$.}
\label{tab:audit-summary}
\end{table}

\subsection{Diagonal mega-shift identity (diagnostic)}
\label{sec:diagonal-megashift}

A useful diagnostic emerged from the branch-gluing audit. For any
residue $r$, the shift $k_\Delta = 2520\, r$ gives
\[
F_{r, k_\Delta}(s)
\;=\; 2520(M s + r) - 2520\, r
\;=\; 2520\, M s
\;=\; A\, s.
\]
Thus a kernel class $s \equiv a \pmod{L}$ can be killed at the
diagonal shift if the class forces sufficiently many prime divisors
into $a$: the forced divisor $D_{k_\Delta}(a, L) = \gcd(A\, a, A\, L)$
inherits all prime contributions of $A$ together with all
contributions of the chosen $a$ within the kernel modulus $L$. This
explains the large ``mega-shift'' kills observed in the audit when
the chosen kernel base happens to be divisible by many primes.

The phenomenon is \emph{not}, by itself, a base-AP closure: Phase B
(below) found alternative clean kernel bases for the sentinel
residues $r \in \{1716, 4862\}$ satisfying the protected
constraints, avoiding the diagonal, and surviving the structural
scan to $10^7$. The diagonal mega-shift is therefore a
\emph{class-specific} structural artifact tied to the particular
kernel base chosen by the SAT solver, not a residue-level closure
mechanism.

\subsection{Branch-gluing Phase B (B1 closure-capable mode)}
\label{sec:branch-gluing-phaseB}

The branch-gluing experiment tested whether independent survivor
kernels at increasing depths were merely incompatible artifacts.
Selected $K = 1021$ branches for $r \in \{1716, 4862\}$ are
diagonal-killed (\cref{sec:diagonal-megashift}), but the
closure-capable B1 search found alternative clean witnesses for
both sentinels. Specifically, under the necessary mode
$\tau(D_k) \le k + 2$ at protected $k \le 1021$ together with
$\tau(D_{k_\Delta}) \le k_\Delta + 2$, with the restriction
$v_p(a) = 0$ for $p \mid A$ and $v_p(a) \le 1$ for cap-2 free primes
$p \in \{23, 29, 31\}$, the adversary CSP returns SAT for both
sentinels. The witnesses satisfy the protected-window constraints,
avoid the diagonal mega-shift, and survive the structural scan to
$10^7$ with the full $172$-prime detector.

Per the standard interpretation (restricted SAT $\Rightarrow$
unrestricted SAT), the branch-gluing mechanism, in the tested
restricted B1 form, does \emph{not} produce a residue closure for
$r \in \{1716, 4862\}$. We do not claim a refutation of all nonlocal
coupling; only that the tested branch-gluing closure mechanism is
negative on the chosen sentinels.

\subsection{Composite-cover audit}
\label{sec:composite-cover}

Finally, we tested whether composite or non-prime-power covering
moduli produce a genuinely new finite-local certificate. The
adversary-SAT \emph{CompositeCoverCert} search asks: does there
exist a finite valuation-pattern cover such that every CRT class
$s \pmod{L}$ is killed by some shift $k \le K_{\max}$? An UNSAT
result would correspond to a finite cover and a closure candidate;
SAT would mean the adversary escapes.

The search returned SAT escapes on four sentinels
$r \in \{0, 1716, 4862, 37752\}$, first at
$(P_{\max}, K_{\max}) = (113, 1000)$ and again at
$(P_{\max}, K_{\max}) = (199, 5000)$. Every generated obstruction
clause reduced to a valuation-pattern clause over prime powers in
the kernel structure. This supports the view that, for the present
framework, composite covering systems do not add closure power
beyond generalized DivisorMask-style valuation covers: the composite
cover encoding is mathematically a generalized DivisorMask at
composite moduli with disjunctive cover semantics, and a SAT escape
strengthens the L1--L3 DivisorMask negative result rather than
competing with it.

\subsection{Small-$k$ quotient-shape compatibility (subsumed)}
\label{sec:small-k-shape-compatibility}

A natural extension of the audit chain considers the
\emph{abstract shape} of the unknown quotient
$Q_k(Y) := F_k(Y) / D_k(a)$ across small shifts $k \le K_{\rm small}$,
subject to the multi-shift difference identity
$F_i(Y) - F_j(Y) = j - i$ and to the per-shift $\tauf$-budget
$\tauf(F_k(Y)) \le k + 2$. The question is whether the abstract
factorization shapes of $Q_1, Q_2, \ldots, Q_{K_{\rm small}}$ can
coexist; an UNSAT result on such an abstract-shape CSP would be a
closure candidate.

The attack is, however, \emph{subsumed by Hensel admissibility}.
For every residue $r \in \Ropen$, the prime-residual Hensel kernel
of \cref{sec:hensel-kernel-audit} satisfies the Conjecture B
condition $2\, \tauf(D_k(a)) \le k + 2$ for all $k \le 1021$. This
directly implies $\tauf(D_k) \le (k+2)/2$, and the abstract
assignment ``$Q_k(Y) = $ a fresh large prime, distinct across
shifts'' satisfies the per-shift $\tauf$-budget:
\[
\tauf(D_k \cdot Q_k) = \tauf(D_k) \cdot \tauf(\text{prime})
= 2 \, \tauf(D_k) \le k + 2.
\]
Fresh large primes with $P > K_{\rm small}$ are automatically
distinct across shifts, since $P \mid F_i$ and $P \mid F_j$ would
force $P \mid (j - i)$, contradicting $|j - i| \le K_{\rm small} < P$.
Multi-shift compatibility is therefore satisfied at the abstract
level by this assignment. We verified this directly: the
fresh-prime assignment produces a valid SAT witness for the
shape-CSP at every $r \in \Ropen$ for every
$K_{\rm small} \in \{16, 32, 64, 1021\}$ (data:
\texttt{data/fresh\_prime\_shape\_witness\_check\_20260425.json}).

The small-$k$ quotient-shape attack therefore adds no closure
power beyond the existing Hensel admissibility data; it confirms
the Schinzel-pattern barrier rather than producing closure.

\subsection{Audit-chain close-out theorem}
\label{sec:audit-chain-closeout}

We collect the negative results assembled in this section into a
single named theorem.

\begin{theorem}[Audit-chain close-out: finite-local forced-divisor
closure-incapability]
\label{thm:audit-chain-closeout}
Within the class of finite-local forced-divisor closure mechanisms
tested in this paper --- comprising:
\begin{enumerate}[leftmargin=*, label=(\roman*)]
\item single-shift forced-divisor primitives (cycle-3 valuation
  leaves and the DivisorMask L1--L3 family;
  \cref{sec:cycle3-audit,sec:divisormask-audit});
\item single-base multi-shift CSP encodings (the prime-residual
  Hensel-balanced kernel framework, the kernel-space CSP, and the
  full safe-set CSP;
  \cref{sec:hensel-kernel-audit,sec:interval-cover-audit});
\item disjunctive multi-shift cover encodings (composite /
  non-prime-power covering certificates, mathematically a
  generalized DivisorMask at composite moduli with disjunctive cover
  semantics; \cref{sec:composite-cover}); and
\item branch-gluing extensions of the single-base CSPs to the
  nested-survivor closure-capable form (Phase B / B1 mode with
  diagonal $\tau(D_{k_\Delta}) \le k_\Delta + 2$ at
  $k_\Delta = 2520\, r$, protected $\tau(D_k) \le k + 2$ at all
  $k \le 1021$, and structural scan to $K_{\rm scan} = 10^7$;
  \cref{sec:branch-gluing-phaseB});
\end{enumerate}
each tested mechanism admits an adversary-SAT survivor configuration
on at least one open residue $r \in \Ropen$, at the parameter
regimes specified in \cref{tab:audit-summary}. Two of the mechanisms
have the strongest possible coverage in this work:
\begin{itemize}[leftmargin=*]
\item the disjunctive multi-shift cover (composite-cover certificate)
  is closure-incapable on \emph{all} 41 residues at
  $P_{\max} = 199$, $K_{\max} = 5000$;
\item the diagonal-B1 mechanism (B1 necessary mode plus diagonal at
  $k_\Delta = 2520\,r$, no scan window) is closure-incapable on $40$ of
  the $41$ residues; the remaining residue $r = 0$ has $k_\Delta = 0$,
  which makes the diagonal constraint $\tau(D_0) \le 2$ degenerate
  (since $\tau(A) = 768 \gg 2$ for any kernel base) and therefore an
  edge case rather than a closure signal.\footnote{The diagonal
  mega-shift identity $F_{r, k_\Delta}(s) = A \cdot s$ requires
  $k_\Delta \ge 1$ to be a valid $\tauf$-barrier shift. At $r = 0$,
  $k_\Delta = 2520 \cdot 0 = 0$ is not a witness shift in the
  $\tauf$-barrier framework. The standard $s \ge 1$ argument applies
  at $r = 0$ in exactly the same way as at the sentinels: the
  $s = 0$ case (corresponding to $n = 2520 \cdot r = 0$, well below
  the $n > 24$ threshold of \cref{prob:e647}) is automatically
  excluded; the diagonal mechanism is parameterized by $s \ge 1$
  for which $n = 2520 \cdot M \cdot s + 2520 \cdot r$ exceeds the
  threshold and the standard analysis applies. The $r = 0$ entry in
  the all-41 triage is therefore an algorithmic artifact at the
  $k_\Delta = 0$ boundary, not an open closure question.}
\end{itemize}
Equivalently: no mechanism in this scope class produces a closure of
the residual frontier $\Ropen$ at the scales tested in this work.

Within this scope class, the audit chain has reached saturation:
every closure-capable mechanism enumerated has been tested at all
$41$ open residues (or, where computational budget permitted only
sentinel coverage, at sentinels chosen for structural
representativeness) and refuted. Closure of the open frontier via
methods within this scope class is therefore ruled out at the
scales tested. Methods outside this scope class --- asymptotic
density, sieve-theoretic pointwise control, conditional analytic
resolution --- are deferred to future work and are addressed
separately by the four convergent obstructions of
\cref{sec:obstructions}.

A diagnostic structural identity
$F_{r, 2520\, r}(s) = A \cdot s$ (\cref{sec:diagonal-megashift})
produces a class-specific kill at $k_\Delta = 2520\, r$ when the
chosen kernel base forces sufficient prime divisors into $a$, but
Phase B exhibits clean alternative kernel bases avoiding the
diagonal for the sentinel residues $r \in \{1716, 4862\}$. The
identity is therefore a uniform structural constraint on admissible
kernel bases, not by itself a closure mechanism.
\end{theorem}

\begin{remark}[Scope of the close-out]
\label{rem:closeout-scope}
\Cref{thm:audit-chain-closeout} concerns the finite-local
forced-divisor scope class enumerated above. It does not address:
\begin{itemize}[leftmargin=*]
\item joint-constraint refinements via cross-shift gcd identities
  $\gcd(F_i(s), F_j(s)) \mid (j - i)$. For any honest CRT class
  $s \equiv a \pmod L$, the identity $F_i(s) - F_j(s) = j - i$ is
  structural, so the resulting pair-level forced-divisor constraint
  $\gcd(D_i, D_j) \mid (j - i)$ is automatically satisfied by
  per-shift valuation constraints.\footnote{Empirically confirmed
  by a pair-subsumption audit on ten witnesses (Phase B SAT
  escapes, composite-cover SAT escapes, and the K=1021
  Hensel-balanced kernel bases): $0$ violations across $\sim 199{,}000$
  pairs per witness, with $\sim 7{,}700$ entangled pairs per witness
  ($\gcd(D_i, D_j) > 1$) all satisfying the divisibility constraint.
  Audit data: \texttt{data/pair\_csp\_subsumption\_audit\_20260425.json}.}
  A non-tautological pair- or multi-shift CSP would need to encode
  joint $\tau$-profile or quotient-shape constraints rather than
  mere forced-divisor compatibility;
\item averaged-sieve methods (Selberg upper bounds,
  Shiu / Nair--Tenenbaum lower bounds, GPY / Maynard /
  Tao--Ter\"av\"ainen multi-point correlation), which operate on
  density rather than pointwise local data; addressed separately by
  the four convergent obstructions of \cref{sec:obstructions};
\item analytic-NT inputs at $q = M$ (effective zero-density
  estimates for $L(s, \chi)$, Conrey--Gonek--Keating asymptotics);
\item methods exploiting global multiplicative or additive structure
  not visible to per-prime valuation analysis.
\end{itemize}
The close-out is thus a precise negative statement within a
precisely-defined scope class, not a refutation of all finite-local
methods.
\end{remark}

\subsection{What the audit does and does not rule out}
\label{sec:audit-scope}

The interpretation of \cref{thm:ple-insufficiency,thm:interval-cover-insufficiency}
requires care.

\begin{remark}[Scope of the audit chain]
\label{rem:audit-scope}
The audit chain rules out the finite-local forced-divisor certificate
families described in this section, at the parameter regimes tested.
Concretely, no \PLEclass\ primitive of the form enumerated in
\cref{def:ple,def:full-safe-set-csp} produces a closure certificate
for any residue $r \in \Ropen$ at the parameter regimes
$K_0 \in \{95, 200, 1021\}$, $B \le 1000$, $K_{\rm scan} \le 10^5$.

\emph{These results do not rule out arbitrary analytic or nonlocal
methods.} They rule out the finite-local forced-divisor certificate
families described in this section. In particular, the audit chain
does not address:
\begin{itemize}[leftmargin=*]
\item averaged-sieve methods of Selberg, Shiu, Nair--Tenenbaum,
  GPY, Maynard, or Tao--Ter\"av\"ainen type, which operate on
  density rather than pointwise local data;
\item Erd\H{o}s covering systems with non-prime-power moduli, when
  the moduli are not subsumed by the CRT decomposition of the
  prime-residual kernel framework;
\item analytic-NT inputs derived from effective zero-density
  estimates for $L(s, \chi)$ at $q = M$, including potential outputs
  of a re-targeted Platt verification;
\item methods that produce a deterministic finite cover of the
  residual frontier rather than a per-shift forced-divisor
  certificate;
\item non-local primitives with a global multiplicative or additive
  structure not visible to per-prime valuation analysis.
\end{itemize}

\Cref{sec:obstructions} addresses the first three of these
explicitly and shows that all are also blocked, but by separate
analytic obstructions rather than by the finite-local audit.
\end{remark}

The combination of the finite-local audit chain
(\cref{sec:audit-summary}) and the analytic obstruction analysis
(\cref{sec:obstructions}) constitutes the barrier of the title:
under any of the standard pointwise-local-existential or
density-averaged sieve approaches, no closure of the residual
frontier follows from the data assembled in this paper.

% =============================================================================
% End of section 6
% =============================================================================
